\chapter{Pattern recognition}
\label{ch:patternrecognition}

Pattern recognition turns the hit cloud of one event (an \code{AtEvent}, produced by the pulse-shape
analysis) into track candidates (\code{AtTrack} objects collected in an \code{AtPatternEvent}). It is the
step between the PSA and the track fit. Besides grouping hits into tracks, the pattern-recognition stage
in ATTPCROOTv2 also prepares everything the fitters consume: it groups the hits of each track into
\code{AtHitCluster} objects with a covariance matrix, orders those clusters, and estimates the initial
geometric parameters of each track (circle centre and radius in the pad plane, the polar angle
\code{GeoTheta}, the azimuth \code{GeoPhi} and a charge sign). Several of the decisions taken here, in
particular the forward/backward direction of a track, are used unchanged by the fitters.

The code lives in \url{AtReconstruction/AtPatternRecognition/} (algorithms) and
\url{AtReconstruction/} (FairRoot tasks). The cluster builder is in \url{AtTools/AtTrackTransformer.cxx}.
Table~\ref{tab:pr:overview} lists the components that are compiled into the framework
(\url{AtReconstruction/CMakeLists.txt}).

\begin{table}[htbp]
\centering
\begin{tabular}{p{0.45\textwidth}p{0.48\textwidth}}
\toprule
Class & Role \\
\midrule
\code{AtPRAtask} & FairTask that runs an \code{AtPRA} algorithm on every event. \\
\code{AtPATTERN::AtPRA} & Abstract base: clustering, cluster ordering, initial parameters,
  primary/fragment selection, kNN pruning. \\
\code{AtPATTERN::AtTrackFinderTC} & TriplClust: triplet building plus hierarchical clustering. \\
\code{AtPATTERN::AtTrackFinderHDBSCAN} & HDBSCAN* density clustering plus cluster joining
  (Spyral style). \\
\code{AtPATTERN::AtTrackFinderRiemann} & Seeded circle RANSAC with a 3D inlier test. \\
\code{SampleConsensus::AtSampleConsensus} & Generic sample consensus (RANSAC, LMedS, MLESAC, \ldots)
  on an \code{AtPattern}. \\
\code{AtSampleConsensusTask}, \code{AtRansacTask} & FairTasks around \code{AtSampleConsensus}. \\
\code{AtPatternBridgeTask} & Fills clusters and \code{Geo*} fields on sample-consensus output. \\
\code{AtTools::AtTrackTransformer} & Builds \code{AtHitCluster}s and their covariance. \\
\code{AtFindVertex} & Vertex finding from straight-line tracks (in \url{AtTools/}). \\
\bottomrule
\end{tabular}
\caption{Pattern-recognition components.}
\label{tab:pr:overview}
\end{table}

The source tree also contains \code{AtTrackFinderHC} (an older PCL-based triplet clustering) and the
directory \url{AtPatternRecognition/trackfinder/}; neither is listed in \url{AtReconstruction/CMakeLists.txt},
so they are not built. There is no Hough-transform track finder: \code{AtPRAtask::SetPRAlgorithm(2)}
only prints ``Using Hough transform algorithm'' (see Section~\ref{sec:pr:task}).

\section{Data flow and the output: \code{AtPatternEvent}}
\label{sec:pr:output}

The input is an \code{AtEvent} branch. In the AT-TPC chains it is \code{AtEventH} (PSA output), a
space-charge corrected copy (\code{AtEventCorrected}) or a cleaned copy (\code{AtEventClean}). The output
is one \code{AtPatternEvent} per event, stored in a \code{TClonesArray} branch (default name
\code{AtPatternEvent}). An \code{AtPatternEvent} holds
\begin{itemize}
\item \code{GetTrackCand()}: a \code{std::vector<AtTrack>} of track candidates;
\item \code{GetNoiseHits()}: the hits that the algorithm did not assign to any track.
\end{itemize}

Each \code{AtTrack} filled by the pattern recognition carries
\begin{itemize}
\item the raw hits (\code{GetHitArray()}, a vector of \code{std::unique\_ptr<AtHit>});
\item the hit clusters (\code{GetHitClusterArray()}, a pointer to a \code{std::vector<AtHitCluster>});
  the fitters (\code{AtGenfitter}, the UKF) read this array, not the raw hits;
\item \code{GetTrackID()}: an integer identifier (see the pitfall in Section~\ref{sec:pr:pitfalls});
\item \code{GetGeoRadius()}, \code{GetGeoCenter()}: circle fitted to the pad-plane projection, in mm;
\item \code{GetGeoTheta()}, \code{GetGeoPhi()}: initial polar and azimuthal angles, in radians;
\item \code{GetChargeSign()}: $+1$ or $-1$ from the rotation sense of the projection, $0$ if unknown;
\item \code{GetPattern()}: the \code{AtPattern} (here an \code{AtPatternCircle2D}) used for the estimate.
\end{itemize}
All positions are in the digitisation (``digi'') frame of the hits: $z$ is the drift coordinate
computed from the time bucket, and the beam travels from high to low $z_\mathrm{digi}$ (comment in
\code{AtPRA::SelectAndMergeTracks}).

\section{The task: \code{AtPRAtask}}
\label{sec:pr:task}

\code{AtPRAtask} (\url{AtReconstruction/AtPRAtask.h}) reads the input \code{AtEvent}, passes it to the
\code{FindTracks} method of an \code{AtPRA} algorithm and writes the returned \code{AtPatternEvent}. It has two
construction modes.

\paragraph{Dependency injection (recommended).} One constructor of \code{AtPRAtask} takes a
\code{std::unique\_ptr<AtPATTERN::AtPRA>}, a fully configured algorithm. All algorithm
settings must then be made on the algorithm object before it is moved into the task; the task does not
forward its own cluster radius, cluster distance or triplet parameters to an injected algorithm. It still
forwards in \code{Init()}: the diffusion parameters from \code{AtDigiPar} and \code{SetPadResXY} (see
Section~\ref{sec:pr:clusters}), and, when enabled on the task, the arc-walk options, the kNN pruning
options, the drift-aware weights and the pre-clustering of the circle fit.

\paragraph{Legacy algorithm number.} With the default constructor, \code{SetPRAlgorithm(int)} selects the
algorithm in \code{Init()}: 0 builds \code{AtTrackFinderTC} (default), 3 builds
\code{AtTrackFinderRiemann}. Values 1 (``RANSAC'') and 2 (``Hough transform'') print a message but
construct no algorithm; \code{Exec()} then calls \code{FindTracks} on an empty pointer. Use
\code{AtSampleConsensusTask} for RANSAC instead (Section~\ref{sec:pr:sac}). HDBSCAN has no algorithm
number and is only available through injection.

\begin{table}[htbp]
\centering
\begin{tabular}{p{0.36\textwidth}p{0.2\textwidth}p{0.36\textwidth}}
\toprule
Setter & Default & Meaning \\
\midrule
\code{SetInputBranch} & \code{AtEventH} & input \code{AtEvent} branch \\
\code{SetOutputBranch} & \code{AtPatternEvent} & output branch \\
\code{SetPersistence} & \code{kFALSE} & write the output branch to file \\
\code{SetMinNumHits}, \code{SetMaxNumHits} & 10, 5000 & event processed only if
  $10 < N_\mathrm{hits} < 5000$ (strict) \\
\code{SetPRAlgorithm} & 0 & legacy selector (0 TC, 3 Riemann) \\
\code{SetScluster}, \code{SetKtriplet}, \ldots & see Sec.~\ref{sec:pr:tc} & TC parameters, legacy path
  only \\
\code{SetClusterRadius}, \code{SetClusterDistance} & 20, 15 mm & cluster builder, legacy path only \\
\code{SetUseArcWalk}, \code{SetTargetClusters}, \code{SetMinHitsPerCluster}, \code{SetArcWalkKNN} &
  false, 25, 3, 10 & arc-walk clustering, forwarded in both modes \\
\code{SetPrunning}, \code{SetkNN}, \code{SetStdDevMulkNN}, \code{SetkNNDist} & off, 5, 0, 10 mm &
  kNN outlier pruning \\
\code{SetPadResXY} & $-1$ (keep 2.3 mm) & transverse pad resolution for the covariance \\
\code{SetUseDriftAwareWeights} & false & weight the circle fit by $1/\sigma^2_{xy}(z)$ \\
\code{SetPreclusterRadiusFit}, \code{SetPreclusterBin} & false, 3 mm & arc-binned refit of the circle \\
\code{SetTCUseSelectAndMerge} & true & TC primary/fragment heuristic (legacy path) \\
\code{SetRiemann*} & see Sec.~\ref{sec:pr:riemann} & Riemann parameters (legacy path) \\
\bottomrule
\end{tabular}
\caption{\code{AtPRAtask} configuration. Defaults from \code{AtPRAtask.h} and the default constructor.}
\label{tab:pr:task}
\end{table}

If \code{FindTracks} returns a null pointer (TC does this for an empty cloud or when the nearest-neighbour
scale \code{dnn} is zero), the task logs a warning and writes no entry for that event: the output
\code{TClonesArray} is empty. Events outside the hit-number window also produce an empty array. Readers
must test \code{GetEntriesFast()} before \code{At(0)}.

\section{The common base: \code{AtPRA}}
\label{sec:pr:base}

\code{AtPATTERN::AtPRA} (\url{AtPatternRecognition/AtPRA.h}) defines the interface
\code{virtual std::unique\_ptr<AtPatternEvent> FindTracks(AtEvent \&event) = 0} and the post-processing
shared by the finders. Which finder calls which step is summarised in Table~\ref{tab:pr:steps}.

\subsection{Initial parameters and the GeoTheta direction}
\label{sec:pr:geotheta}

\code{SetTrackInitialParameters(AtTrack\&)} (protected) computes the \code{Geo*} fields:
\begin{enumerate}
\item A subset of the track's raw hits is taken from the end of the hit array, which the code comment
  calls the vertex end (``highest Z\_digi = last in array''). Its size is
  $\max(\min(n_\mathrm{min},N), \min(\max(3, fN), n_\mathrm{max}, N))$ with $f$ =
  \code{SetRadiusFitFraction} (default 1.0, i.e.\ all hits), $n_\mathrm{min}$ = \code{SetMinHitsRadius}
  (3) and $n_\mathrm{max}$ = \code{SetMaxHitsRadius} (1000).
\item A circle RANSAC (\code{AtSampleConsensus}, pattern \code{kCircle2D}, distance threshold 6~mm,
  1000 iterations, minimum $0.1\times$ the subset size; all hard-coded) selects the inliers.
\item An \code{AtPatternCircle2D} is fitted to the whole subset; optionally re-fitted with drift-aware
  weights (\code{SetUseDriftAwareWeights}) or on charge-weighted centroids of arc bins of width
  \code{SetPreclusterBin} (\code{SetPreclusterRadiusFit}). Centre and radius become \code{GeoCenter} and
  \code{GeoRadius}.
\item The charge sign is taken from the sign of the cross product of the hits ordered by distance from the
  beam axis: \code{crossZ > 0} gives $-1$, otherwise $+1$. The code comment states that this was
  calibrated on PUMA simulation.
\item For the circle inliers, in hit-array order, the azimuth around the centre is unwrapped by counting
  crossings of the $-y$ axis, and each hit is mapped to the plane (arc length $s$, $z$), with $s$ measured
  from the first inlier.
\item A line RANSAC (pattern \code{kLine}, threshold 6~mm, minimum $0.1\times$ the number of points, line
  refit on) is run in the $(s,z)$ plane. Only the first line is used, i.e.\ the one with most inliers
  (the RANSAC estimator stores $\chi^2 = 1/N_\mathrm{inliers}$ and patterns are sorted by it). With
  $\hat d=(d_s,d_z)$ the unit direction of that line, the code sets
  \[
  \theta_\mathrm{geo} = \arccos\bigl(\mathrm{sign}\cdot|d_z|\bigr),\qquad
  \mathrm{sign} = -1 \text{ if } d_s d_z < 0,\ +1 \text{ otherwise.}
  \]
  \code{GeoPhi} is the azimuth of the first inlier around the circle centre.
\end{enumerate}

The forward/backward decision is therefore the sign of the slope $\mathrm{d}z/\mathrm{d}s$ of a single
RANSAC line: $\theta_\mathrm{geo} > 90^\circ$ when the slope is negative. Downstream code relies on it:
\code{AtGenfitter} treats a track as backward when \code{GetGeoTheta() > M\_PI/2}
(\url{AtFitter/AtGenfitter.cxx}), and \code{AtGenfit} chooses its forward or backward branch from the same angle (with $180^\circ -
\theta_\mathrm{geo}$ under its simulation convention). Because the arc length is built by unwrapping in
hit-array order, the result depends on that order (Section~\ref{sec:pr:pitfalls}). The RANSAC samples with
\code{gRandom} (\url{AtTools/AtHitSampling/AtUniform.cxx}), so the result also depends on the state of the
random generator. If no line is found the angle stays at its initial value 0, and the track is treated as
forward. The macro \url{macro/Unpack_HDF5/a1975/D2_UKF/pra_theta.C} replays this calculation for one track
and draws the $(s,z)$ plane. \code{AtGenfitter} offers alternative sources for the direction
(\code{SetDirectionFromSpyral}, \code{SetDirectionFrom3D}); these are described in the fitting chapter.

\subsection{Cluster ordering}

\code{OrderClustersAlongTrack(AtTrack\&)} reorders the cluster array by a greedy nearest-neighbour walk
starting from the cluster with the highest $z$ (the vertex end for forward tracks). The walk must consume
every cluster; clusters it skips are visited at the end, so the tail of the array can jump back to the
middle of the track. \url{macro/Unpack_HDF5/a1975/D2_UKF/zuturn_dp.C} measures this ``U-turn'' (total
variation of $z$ along the array divided by the $z$ span) and quotes run\_0016 event 31316 as an example.
The function reorders clusters only; the raw hit array is untouched.

\subsection{Primary selection and fragment merging}
\label{sec:pr:selectmerge}

\code{SelectAndMergeTracks(tracks, vertexRadiusXY = 80, mergeDist = 30, minLabTheta = 10)} is an AT-TPC
heuristic used by TC and Riemann (switchable, on by default). The finders call it with the defaults, so
the three arguments cannot be changed from a macro. On the cluster arrays it
\begin{enumerate}
\item defines a vertex end (cluster closest to the beam axis in $xy$) and a far end (farthest from the
  axis) for each track, independent of $z$;
\item rejects tracks whose vertex-to-far-end direction is within 10$^\circ$ of the beam axis
  ($\theta_\mathrm{lab}<10^\circ$ or $>170^\circ$) and degenerate tracks with no extent;
\item marks as primary the tracks whose vertex end is within 80~mm of the beam axis and within 50~mm
  (hard-coded) in $z$ of the highest vertex-end $z$ of the event;
\item appends to a primary the hits of any non-primary track with an end within 30~mm of the primary's far
  end, re-clusters the merged track and reorders its clusters; this repeats until nothing merges;
\item deletes every remaining non-primary track.
\end{enumerate}
Two consequences follow from the code. Hits of rejected tracks are not added to the noise of the
\code{AtPatternEvent}; they disappear from the output. A track whose clusters are empty (for example
because the cluster radius is 0, Section~\ref{sec:pr:clusters}) has a degenerate direction and is rejected
in step 2. Disable the heuristic with \code{SetUseSelectAndMerge(false)} on TC or Riemann for annular geometries
such as PUMA, as the header comments advise. The 50~mm window assumes that the primaries of an event
share one vertex.

\subsection{kNN pruning}

With \code{SetPrunning()} each track is passed to \code{PruneTrack}: for every hit the mean $\mu$ and
standard deviation $\sigma$ of the distances to its \code{SetkNN} (5) nearest hits (the hit itself counts, at distance 0) are computed, and the
hit is erased if $\mu + \sigma\cdot$\code{SetStdDevMulkNN} (0) $\geq$ \code{SetkNNDist} (10~mm). In TC
and Riemann pruning runs after the clusters are built, so pruned hits still contribute to the clusters.
HDBSCAN does not call it.

\subsection{Unused members}

\code{SetMaxHits}, \code{SetMinHits}, \code{SetMeanDistance} and \code{SetUseHitsForRadius} store values
that no finder reads in the current code; \code{AtTrackFinderTC::SetPadding} likewise stores a value that
is not passed to the triplet clustering.

\begin{table}[htbp]
\centering
\begin{tabular}{p{0.33\textwidth}p{0.17\textwidth}p{0.17\textwidth}p{0.2\textwidth}}
\toprule
Step & TC & HDBSCAN & Riemann \\
\midrule
cluster builder & Smooth3D or arc walk & Smooth3D only & Smooth3D or arc walk \\
default radius / distance & 0 / 0 (must be set) & 15 / 7.5 mm if unset & 0 / 0 (must be set) \\
\code{OrderClustersAlongTrack} & yes & yes & yes \\
\code{SelectAndMergeTracks} & yes (switch) & no & yes (switch) \\
\code{SetTrackInitialParameters} & after selection & before clustering & after selection \\
drops NaN/$\le$0 \code{GeoRadius} & yes & no & no \\
kNN pruning & optional & no & optional \\
track ID & cluster label & 0, 1, 2, \ldots & order found \\
\bottomrule
\end{tabular}
\caption{Post-processing performed by each finder (from the \code{FindTracks} implementations).}
\label{tab:pr:steps}
\end{table}

\section{TriplClust: \code{AtTrackFinderTC}}
\label{sec:pr:tc}

\paragraph{Idea.} \code{AtTrackFinderTC} wraps the TriplClust algorithm of Dalitz, Wilberg and Aymans
(``TriplClust: An Algorithm for Curve Detection in 3D Point Clouds'', IPOL; source in
\url{AtPatternRecognition/triplclust/}). A track is a curve, locally a straight line. The algorithm
\begin{enumerate}
\item computes $d_\mathrm{NN}$, the square root of the first quartile of the squared nearest-neighbour
  distances of the cloud;
\item smooths the cloud by replacing each point with the mean of its neighbours within radius $r$;
\item for each smoothed point, forms triplets with pairs of its $k$ nearest neighbours whose two branches
  are nearly collinear ($1-\cos\angle \le a$) and keeps the $n$ best;
\item clusters the triplets by single-linkage hierarchical clustering with a triplet distance scaled by $s$,
  and cuts the dendrogram at distance $t$;
\item discards clusters with fewer than $m$ triplets and converts triplet clusters into point clusters.
\end{enumerate}
Each surviving cluster becomes one \code{AtTrack}; points in no cluster become noise. A point that belongs
to triplets of two clusters is added to both tracks.

\paragraph{Configuration.} The setters act on \code{inputParams} (defaults in
\code{AtTrackFinderTC.h}, identical to the legacy defaults in the \code{AtPRAtask} constructor). In the
wrapper $r$ and $s$ are multiples of $d_\mathrm{NN}$ (\code{rdnn} and \code{sdnn} are true in
\code{Opt::Opt()}), and the automatic cut on $t$ is off.

\begin{table}[htbp]
\centering
\begin{tabular}{p{0.32\textwidth}p{0.1\textwidth}p{0.5\textwidth}}
\toprule
Setter & Default & Meaning \\
\midrule
\code{SetRsmooth} & 2 & smoothing radius, in units of $d_\mathrm{NN}$ \\
\code{SetKtriplet} & 19 & neighbours considered when building triplets \\
\code{SetNtriplet} & 2 & best triplets kept per point \\
\code{SetAtriplet} & 0.03 & maximum $1-\cos$ of the angle between the triplet branches \\
\code{SetScluster} & 0.3 & distance scale of the triplet metric, in units of $d_\mathrm{NN}$ \\
\code{SetTcluster} & 4.0 & dendrogram cut distance \\
\code{SetMcluster} & 15 & minimum number of triplets per cluster \\
\code{SetUseSelectAndMerge} & true & apply Section~\ref{sec:pr:selectmerge} \\
\code{SetClusterRadius}, \code{SetClusterDistance} & 0, 0 & cluster builder (base-class defaults; must
  be set when injecting) \\
\bottomrule
\end{tabular}
\caption{\code{AtTrackFinderTC} parameters.}
\end{table}

After clustering, each track is clustered (\code{ClusterizeSmooth3D} with the base-class radius and
distance, or \code{ClusterizeArcWalk}), its clusters are ordered, optionally pruned, then
\code{SelectAndMergeTracks} runs, then \code{SetTrackInitialParameters}. Tracks whose \code{GeoRadius} is
NaN or $\le 0$ are dropped with a message ``AtPRA filter: \ldots dropped''; their hits are not put into
the noise.

\paragraph{Example.} The a1975 macro \url{macro/Unpack_HDF5/a1975/UKF/pipeline/unpackReco_a1975_UKF.C}
(16C on H$_2$) uses TC with the default triplet parameters and sets only the cluster builder:
\begin{verbatim}
   Double_t clusterRadius = 15.0;
   Double_t clusterDistance = 7.5;
   auto praAlgo = std::make_unique<AtPATTERN::AtTrackFinderTC>();
   praAlgo->SetClusterRadius(clusterRadius);
   praAlgo->SetClusterDistance(clusterDistance);
   AtPRAtask *praTask = new AtPRAtask(std::move(praAlgo));
   praTask->SetInputBranch(useSC ? "AtEventCorrected" : "AtEventH");
   praTask->SetOutputBranch("AtPatternEvent");
   praTask->SetPersistence(true);
\end{verbatim}
The same pattern appears in \url{macro/Unpack_HDF5/a1975/unpackNFit_a1975.C} and
\url{macro/Unpack_HDF5/a1975/D2_UKF/unpackReco_a1975_deuterium.C}. Other a1975 macros use HDBSCAN
instead (\url{macro/Unpack_HDF5/a1975/D2_UKF/unpackReco_multifit.C},
\url{macro/Unpack_HDF5/a1975/UKF/pd/prod_h2_hdb.sh}), so check the macro of the production in question.

\section{HDBSCAN: \code{AtTrackFinderHDBSCAN}}
\label{sec:pr:hdbscan}

\paragraph{Idea.} \code{AtTrackFinderHDBSCAN} is a self-contained C++ implementation of HDBSCAN* followed by
the cluster joining and outlier cleaning of the Spyral analysis package (header comment: ``Mirrors Spyral
a1975 config''). The header states its purpose: to separate overlapping or crossing tracks that TC merges,
and then to join over-segmented pieces of one track. The steps in \code{FindTracks} are:
\begin{enumerate}
\item \textbf{Core distances.} For each hit, the distance to its $k$-th nearest neighbour, $k$ =
  \code{SetMinSamples}.
\item \textbf{Mutual-reachability MST.} The edge weight between hits $i$ and $j$ is
  $\max(\mathrm{core}_i, \mathrm{core}_j, d_{ij})$; the minimum spanning tree is built by Prim's algorithm over all pairs
  ($O(N^2)$ in time).
\item \textbf{Condensed tree and excess of mass.} Splits that leave fewer than the minimum cluster size
  points are treated as points falling out; the most stable clusters are selected. With
  \code{SetClusterSelectionEpsilon} $>0$, a selected cluster is replaced by its parent while the split
  that created it happened at a distance below epsilon, so pieces closer than epsilon (mm) are not separated.
\item \textbf{Joining} of the resulting clusters by the method chosen with \code{SetJoinMethod}.
\item \textbf{Local Outlier Factor cleaning} inside each cluster: a hit is removed if its LOF exceeds
  \code{SetOutlierThreshold}, with $k=\max(2, \mathrm{scale}\cdot n)$ neighbours.
\item Each cluster becomes an \code{AtTrack} with IDs 0, 1, 2, \ldots; \code{SetTrackInitialParameters}
  is called, then the hit clusters are built with \code{ClusterizeSmooth3D} and ordered. Unassigned hits
  go to the noise.
\end{enumerate}
The minimum cluster size is \code{SetMinClusterSize} if positive, otherwise
$\max(\mathrm{lowerCutoff}, \mathrm{scale}\cdot N_\mathrm{hits})$.

\paragraph{Join methods.} All methods only consider clusters with at least \code{SetMinClusterSizeJoin}
hits.
\begin{itemize}
\item \code{"overlap"} (default): each cluster gets a least-squares (Kasa) circle and a winding direction
  (fraction of increasing azimuth steps along $z$ compared with \code{SetDirectionThreshold}); clusters with
  the same direction merge if the overlap area of their circles exceeds \code{SetCircleOverlapRatio} times
  the area of the smaller circle. Circles with radius below 10~mm are ignored. Iterates until stable.
\item \code{"continuity"}: same-direction clusters merge if the downstream end of one and the upstream end
  of the other are close in $z$ (fraction \code{SetJoinZFraction} of the total $z$ extent) and in radial
  distance from the circle centre (fraction \code{SetJoinRadiusFraction} of the mean radius).
\item \code{"both"}: \code{"continuity"} then \code{"overlap"}.
\item \code{"motion"}: a vertex proxy is computed (centroid of the 10\,\% of clustered hits closest to the
  beam axis); starting from the cluster nearest to it, clusters are chained end to end when the gap is below
  \code{SetMotionGapTol} (mm) and the travel tangents (and the gap direction) agree within
  \code{SetMotionAngleTol} (degrees). This allows the curvature to change along the track, unlike the
  fixed-circle joins.
\item \code{"mover"}: \code{"motion"} followed by \code{"overlap"} (the comment: end-to-end energy loss,
  then spiral splits).
\item \code{"none"}: no joining.
\end{itemize}

\begin{table}[htbp]
\centering
\begin{tabular}{p{0.34\textwidth}p{0.13\textwidth}p{0.45\textwidth}}
\toprule
Setter & Default & Meaning \\
\midrule
\code{SetMinClusterSize} & 0 & fixed minimum cluster size; $\le 0$ selects the adaptive value \\
\code{SetMinSizeScaleFactor} & 0.05 & adaptive size $=\max(\mathrm{cutoff}, 0.05\,N_\mathrm{hits})$ \\
\code{SetMinSizeLowerCutoff} & 10 & lower bound of the adaptive size \\
\code{SetMinSamples} & 3 & $k$ of the core distance \\
\code{SetClusterSelectionEpsilon} & 10.0 & do not split clusters closer than this (mm) \\
\code{SetAllowSingleCluster} & true & allow the root to be the only cluster \\
\code{SetJoinMethod} & \code{"overlap"} & see the list above \\
\code{SetCircleOverlapRatio} & 0.25 & overlap join threshold \\
\code{SetMinClusterSizeJoin} & 15 & minimum hits for a cluster to take part in a join \\
\code{SetJoinZFraction} & 0.2 & continuity join, $z$ gap fraction \\
\code{SetJoinRadiusFraction} & 0.3 & continuity join, radius fraction \\
\code{SetDirectionThreshold} & 0.5 & winding-direction decision \\
\code{SetMotionGapTol} & 40.0 & motion join, maximum gap (mm) \\
\code{SetMotionAngleTol} & 35.0 & motion join, maximum tangent mismatch (deg) \\
\code{SetOutlierScaleFactor} & 0.05 & LOF neighbours fraction; $\le 0$ disables LOF \\
\code{SetOutlierThreshold} & 1.5 & LOF cut \\
\code{SetClusterRadius}, \code{SetClusterDistance} & 15, 7.5 & used when the base-class value is 0 \\
\bottomrule
\end{tabular}
\caption{\code{AtTrackFinderHDBSCAN} parameters (\code{AtTrackFinderHDBSCAN.h} and \code{.cxx}).}
\end{table}

HDBSCAN does not run \code{SelectAndMergeTracks}, does not drop tracks with an invalid radius, ignores
\code{SetUseArcWalk} and does not prune. Beam-like clusters are therefore kept as track candidates and
must be removed downstream.

\paragraph{Example.} The a1975 16C(d,t) simulation reconstruction,
\url{macro/Simulation/ATTPC/16C_dt/run_reco_C16dt.C}, with the macro defaults \code{hdMcs = 20},
\code{hdMs = 8}, \code{joinMethod = "mover"}:
\begin{verbatim}
   auto hdb = std::make_unique<AtPATTERN::AtTrackFinderHDBSCAN>();
   hdb->SetMinClusterSize(hdMcs);
   hdb->SetMinSamples(hdMs);
   hdb->SetClusterSelectionEpsilon(10.0);
   hdb->SetJoinMethod(joinMethod.Data());
   hdb->SetMinClusterSizeJoin(15);
   hdb->SetCircleOverlapRatio(0.25);
   hdb->SetMotionGapTol(40);
   hdb->SetMotionAngleTol(35);
   AtPRAtask *praTask = new AtPRAtask(std::move(hdb));
   praTask->SetInputBranch("AtEventClean");
   praTask->SetOutputBranch("AtPatternEvent");
   praTask->SetPersistence(kTRUE);
\end{verbatim}
No cluster radius is set, so the 15/7.5~mm fallback applies. The header comment of this macro records
that the (p,d) simulation was reconstructed with \code{AtPSAMultiFit} + HDBSCAN while the (p,d) data ran
\code{AtPSAMax} + TC; when an acceptance is computed from simulation, the finder (and the PSA) should match
the one used on the data.

\section{Seeded circle RANSAC: \code{AtTrackFinderRiemann}}
\label{sec:pr:riemann}

Despite its name (kept for the legacy algorithm number 3), the current \code{AtTrackFinderRiemann} is a
circle RANSAC with local seeds, as its header explains. In each of at most \code{SetMaxTracks} passes over
the unassigned hits it draws a random seed hit and two more from its \code{SetKNN} nearest neighbours in 3D,
builds the circle through the three and a line $z = a + b\phi$ through two of them, and counts hits that are
within \code{SetInlierDist} of the circle in $xy$ and within \code{SetZInlierDist} of the line in $z$. The
best candidate is refitted, cut to its largest contiguous arc (gap below \code{SetMaxPhiGap}), optionally
extended by a sliding-window arc walk (\code{SetUseArcWalkExtend}), and accepted with at least
\code{SetMinHitsPerTrack} hits. The random generator is a \code{std::mt19937} seeded with
\code{SetSeed} (12345), so results are reproducible.

\begin{table}[htbp]
\centering
\begin{tabular}{p{0.3\textwidth}p{0.12\textwidth}p{0.5\textwidth}}
\toprule
Setter (task setter) & Default & Meaning \\
\midrule
\code{SetInlierDist} (\code{SetRiemannInlierDist}) & 5.0 (15.0) & $xy$ residual cut, mm \\
\code{SetZInlierDist} & 15.0 & $z$ residual cut, mm \\
\code{SetMinHitsPerTrack} (\code{SetRiemannMinHits}) & 10 & minimum hits per track \\
\code{SetMaxTracks} (\code{SetRiemannMaxTracks}) & 6 & maximum tracks extracted \\
\code{SetMaxIterations} (\code{SetRiemannMaxIter}) & 400 & RANSAC iterations per pass \\
\code{SetKNN}, \code{SetSeed}, \code{SetMaxPhiGap} & 10, 12345, 60$^\circ$ & seeding and arc split \\
\code{SetUseSelectAndMerge} & true & Section~\ref{sec:pr:selectmerge} \\
\code{SetUseArcWalkExtend}, \code{SetArcWalkWindow}, \code{SetArcWalkMaxMiss} & false, 10, 5 & tail
  recovery \\
\bottomrule
\end{tabular}
\caption{\code{AtTrackFinderRiemann} parameters. Values in parentheses are the \code{AtPRAtask}
defaults used on the legacy path, which differ for the inlier distance.}
\end{table}

\section{Sample consensus: \code{AtSampleConsensus} and its tasks}
\label{sec:pr:sac}

\code{SampleConsensus::AtSampleConsensus} (\url{AtPatternRecognition/AtSampleConsensus.h}) fits a
parametric \code{AtPattern} to a hit cloud. It is used in two ways: as a track finder for straight or
simple topologies (fission, beam plus products), and internally by \code{AtPRA} for the circle and the
$(s,z)$ line of Section~\ref{sec:pr:geotheta}.

\paragraph{Algorithm.} \code{Solve} generates \code{SetNumIterations} candidate patterns from random
samples of the minimum number of points of the pattern, scores each with the chosen estimator and keeps
those with more than \code{SetMinHitsPattern} inliers, sorted by the estimator's $\chi^2$ (for RANSAC
$1/N_\mathrm{inliers}$). Then, best first, it moves the hits within \code{SetDistanceThreshold} of each
pattern out of the cloud into a new track (optionally refitting the pattern, \code{SetFitPattern}, with a
charge-weighted fit above \code{SetChargeThreshold}, $-1$ disabling the weighting). Remaining hits are
noise. Track IDs are the positions in the output vector.

\begin{itemize}
\item Patterns (\code{AtPatterns::PatternType}): \code{kLine}, \code{kRay}, \code{kCircle2D}, \code{kY},
  \code{kFission}.
\item Estimators (\code{SampleConsensus::Estimators}): \code{kRANSAC} (inlier count), \code{kLMedS},
  \code{kMLESAC}, \code{kWRANSAC} (charge weighted), \code{kChi2} (mean squared residual of inliers),
  \code{kYRANSAC} (ignores the beam branch of a Y pattern).
\item Samplers (\code{RandomSample::SampleMethod}): \code{kUniform}, \code{kChargeWeighted},
  \code{kGaussian}, \code{kWeightedGaussian}, \code{kWeightedY}, \code{kWeightedGaussianTrunc}, \code{kY}.
\end{itemize}
Defaults of the class: estimator \code{kRANSAC}, pattern \code{kLine}, sampler \code{kUniform}, 500
iterations, 30 minimum points, 15~mm threshold, pattern refit on, no charge weighting.

\paragraph{Tasks.} \code{AtSampleConsensusTask} takes a configured \code{AtSampleConsensus} as a
\code{unique\_ptr}; branches default to \code{AtEventH} in and \code{AtPatternEvent} out, persistence
off. \code{AtRansacTask} builds a line fit itself: estimator \code{kChi2} by default, \code{kMLESAC},
\code{kLMedS} or \code{kWRANSAC} for \code{SetAlgorithm(2)}, \code{(3)}, \code{(4)}; threshold
\code{SetDistanceThreshold} (5~mm), \code{SetMinHitsLine} (5), \code{SetNumItera} (500), sampler
\code{SetRanSamMode} (0, uniform). The \code{fRANSACModel} member set by \code{SetModelType} is not used in
\code{Exec}.

Sample-consensus tracks have no hit clusters and zero \code{Geo*} fields. \code{AtPatternBridgeTask}
(\url{AtReconstruction/AtPatternBridgeTask.h}) fills them for tracks whose pattern is an
\code{AtPatternCircle2D}: it copies centre and radius, runs \code{ClusterizeSmooth3D} with its own
defaults (radius 8~mm, distance 5~mm) and estimates \code{GeoTheta} from the first and last cluster.

\paragraph{Example.} A fission (two-line) fit in \url{macro/e12014/adam/unpack/unpack_linked.C}:
\begin{verbatim}
   auto method = std::make_unique<SampleConsensus::AtSampleConsensus>(
      SampleConsensus::Estimators::kYRANSAC, AtPatterns::PatternType::kFission,
      RandomSample::SampleMethod::kY);
   method->SetDistanceThreshold(20);
   method->SetNumIterations(500);
   method->SetMinHitsPattern(150);
   method->SetChargeThreshold(10); //-1 implies no charge-weighted fitting
   method->SetFitPattern(true);
   auto sacTask = new AtSampleConsensusTask(std::move(method));
\end{verbatim}

\section{Hit clusters for the fit}
\label{sec:pr:clusters}

The fitters do not fit raw hits. Each track's hits are condensed into \code{AtHitCluster}s
(\url{AtData/AtHitCluster.h}, derived from \code{AtHit}) that carry a position, a charge, a mean time
stamp and a $3\times3$ covariance matrix (\code{GetCovMatrix()}). Two builders exist in
\code{AtTools::AtTrackTransformer}.

\subsection{\code{ClusterizeSmooth3D(track, radius, distance)}}

\begin{enumerate}
\item Walk the hit array in its stored order, keeping a reference position (initially the first hit).
  Each time a hit lies at least \code{distance} from the reference, collect all hits of the track within
  \code{radius} of the reference and form a cluster, unless an existing cluster lies within
  \code{distance}; then move the reference to the current hit.
\item Smoothing pass: if there are more than two clusters, for each consecutive pair insert the midpoint,
  and re-cluster around each original centre and each midpoint with radius/2; these clusters replace the
  first-pass ones.
\end{enumerate}
The cluster $x$ and $y$ are charge-weighted means; $z$ and the time stamp are unweighted means. Both
parameters are strict inequalities on 3D distances in mm: with \code{radius} = 0 no hit is ever collected
and the track has no clusters. This is why \code{SetClusterRadius} must be called on an injected TC or
Riemann finder, whose base-class default is 0. The \code{AtPRAtask} legacy defaults are radius 20~mm,
distance 15~mm (comment: ``Gives 2.2\% RMS (was r10 d20 $\to$ 3.8\%)''); the a1975 TC macros use 15 and
7.5~mm, which is also the HDBSCAN fallback.

\subsection{\code{ClusterizeArcWalk(track, targetClusters, minHitsPerCluster, kNN)}}

Enabled with \code{SetUseArcWalk(true)} (TC and Riemann only). It builds a kNN graph of the hits, extracts a
spanning tree and an ordering along it by a double depth-first search, and groups consecutive hits into
clusters of $\max(\mathrm{minHitsPerCluster}, N_\mathrm{hits}/\mathrm{targetClusters})$ hits (defaults 25
target clusters, 3 hits minimum, $k=10$). The ordering does not sort on $z$, which is why the
\code{AtGenfitter} header recommends it together with \code{SetUseClusterOrder}. Note that TC and Riemann
still call \code{OrderClustersAlongTrack} after it, so the stored order is the nearest-neighbour walk from the
highest-$z$ cluster rather than the arc-walk order.

\subsection{Cluster covariance}

Both builders compute a diagonal covariance. For a hit at drift distance $z$ (mm), with drift time
$t = z/(10\,v_d)$ ($v_d$ in cm/\textmu s),
\begin{align*}
\sigma^2_{x} = \sigma^2_{y} &= \sigma^2_\mathrm{pad} + 200\,D_T\,t, \\
\sigma^2_{z} &= (1.5\,\sigma_\mathrm{pad})^2 + \frac{(10\,v_d\,t_\mathrm{TB})^2}{12} + 200\,D_L\,t,
\end{align*}
in mm$^2$, with $D_T$, $D_L$ in cm$^2$/\textmu s and $t_\mathrm{TB}$ the time-bucket length in \textmu s.
The cluster variance is that of a charge-weighted mean, $\sigma^2_\mathrm{cl} = \sum q_i^2\sigma_i^2 /
(\sum q_i)^2$. The defaults in \code{AtTrackTransformer.h} are $D_T = 9\times10^{-5}$, $D_L =
9\times10^{-7}$, $v_d = 1.0$, $t_\mathrm{TB} = 0.320$ and $\sigma_\mathrm{pad} = 2.3$~mm ($8/\sqrt{12}$).
\code{AtPRAtask::Init()} overwrites the first four with \code{AtDigiPar} values
(\code{GetCoefDiffusionTrans}, \code{GetCoefDiffusionLong}, \code{GetDriftVelocity},
\code{GetTBTime}$\times10^{-3}$) and the pad resolution with \code{SetPadResXY} if positive.

\url{docs/development/ClusterCovariance.md} describes the earlier state, in which these coefficients were
hard-coded 10 to 1000 times too large, and the plan (pass \code{AtDigiPar}, use $8/\sqrt{12}$ for the
pad resolution) that the code above now implements. Its Step 4 (validation of $\chi^2/\mathrm{ndf}$ with
the UKF) is not documented as done. \url{docs/development/ClusteringOptimization.md} lists the open
concerns about the cluster builder: midpoint insertion can move cluster centres off the track, the
ordering depends on finding the vertex end, and \code{GeoRadius} is a projection-plane estimate.

\section{Vertex finding: \code{AtFindVertex}}

\code{AtFindVertex} (\url{AtTools/AtFindVertex.h}; a \code{TObject}-derived copy is in
\url{AtAnalysis/}) finds vertices from straight-line tracks: it casts each track's pattern to
\code{AtPatternLine}, so it needs sample-consensus line tracks, not the circle patterns of the helical
finders. It is constructed with a line-distance
threshold (default 15~mm), takes the beam line with \code{SetBeam(pos, dir)}, and
\code{FindVertex(tracks, nbTracksPerVtx)} groups lines whose mutual distance is below the threshold and
provides, through \code{GetTracksVertex()}, a list of vertex positions with their tracks. It is used by the
d2He analysis macros (\url{macro/Simulation/ATTPC/d2He/ana_d2He.C}, with threshold 12~mm). It is not used
by the helical finders above, which locate the vertex only through the heuristics of
Section~\ref{sec:pr:selectmerge}.

\section{Choosing an algorithm}

\begin{itemize}
\item \textbf{TC} is the default of \code{AtPRAtask}; the header of
  \url{macro/Simulation/ATTPC/16C_dt/run_reco_C16dt.C} records it as the finder of the a1975 (p,d) data. With \code{SelectAndMergeTracks} it returns only tracks that start near the beam axis
  at the event vertex, with beam tracks and isolated fragments removed. Its memory grows with the square
  of the number of triplets (condensed distance matrix in \code{compute\_hc}).
\item \textbf{HDBSCAN} is intended for events where tracks overlap or cross, which TC merges (header
  comment), and reproduces the Spyral clustering. Its output keeps every cluster, including beam-like
  ones, and its cost grows as $N_\mathrm{hits}^2$ (all-pairs core distances and MST). Use
  \code{"mover"} or \code{"motion"} joining for stopping particles whose curvature changes.
\item \textbf{Riemann} (seeded circle RANSAC) is deterministic for a fixed seed and finds a fixed maximum
  number of helices; the header documents it as replacing earlier variants that lost too many hits.
\item \textbf{Sample consensus} is the choice for straight tracks without a magnetic field or for
  parametric topologies (\code{kY}, \code{kFission}); add \code{AtPatternBridgeTask} if a helix fitter
  must consume circle patterns.
\end{itemize}
When an acceptance or efficiency is taken from simulation, the simulation must be reconstructed with the
same finder and settings as the data (Section~\ref{sec:pr:hdbscan}).

\section{Known pitfalls}
\label{sec:pr:pitfalls}

\paragraph{Track IDs are not array positions.} In \code{AtTrackFinderTC}, \code{SetTrackID(cluster\_index)}
stores the index of the TriplClust cluster. \code{SelectAndMergeTracks} then erases tracks and the
\code{GeoRadius} filter skips more, so the surviving IDs are in general non-contiguous and
\code{GetTrackCand()[id]} is not the track with that ID. The same holds for Riemann (ID = number of tracks
before it, then selection erases some). Only HDBSCAN and \code{AtSampleConsensus} number tracks by their
position. Loop over \code{GetTrackCand()} and compare \code{GetTrackID()}, never index by it.

\paragraph{Hit order is time order, and it is not an order along the track.} \code{AtPSA::Analyze} ends with
\code{event->SortHitArrayTime()}, a \code{std::sort} on the integer time stamp. \code{std::sort} is not
stable, so hits in the same time bucket have no defined order. Cleaning and space-charge tasks copy hits in input order. TC fills each track in increasing point
index (sorted in \code{cluster\_triplets\_to\_points}), i.e.\ in this time order; HDBSCAN fills each label
in increasing hit index, so a joined track is a concatenation of time-ordered pieces; \code{SelectAndMergeTracks} appends
merged fragments at the end. Three consumers depend on that order:
\begin{itemize}
\item \code{SetTrackInitialParameters} takes the fit subset from the end of the array and unwraps the
  azimuth in array order, so \code{GeoTheta} (and hence forward/backward) depends on it. The header of
  \url{macro/Unpack_HDF5/a1975/D2_UKF/theta_order_test.C} reports that, measured on run\_0020, the stored hit
  path of HDBSCAN tracks is up to 2.8 times the nearest-neighbour path. \code{OrderClustersAlongTrack} does
  not help: it reorders clusters, not hits.
\item \code{ClusterizeSmooth3D} walks the hit array in its stored order to place the cluster centres.
\item The fitters. Unless \code{SetUseClusterOrder(true)} is set, \code{AtGenfitter} sorts the clusters
  by $z$ before fitting. Its header explains why that fails for tracks nearly perpendicular
  to the beam: the $z$ advance per cluster falls below the $z$ noise and one time bucket, so the sort orders
  noise. The option keeps the pattern-recognition order instead; the header pairs it with
  \code{AtPRA::SetUseArcWalk(true)}.
\end{itemize}

\paragraph{GeoTheta is one sign of one RANSAC line.} The direction rests on the sign of the slope of the
best line in the $(s,z)$ plane (Section~\ref{sec:pr:geotheta}). It is random (\code{gRandom}), sensitive
to hit order, defaults to $0$ (forward) when no line is found, and is used as is by \code{AtGenfitter}
(backward if $>90^\circ$) and \code{AtGenfit}. The \code{AtGenfitter} header states that on a1975 H$_2$
data GeoTheta disagrees with the Spyral estimate on 8\,\% of tracks and is the wrong one where the elastic
locus can decide; see \code{SetDirectionFromSpyral} and \code{SetDirectionFrom3D} there.

\paragraph{Merged candidates.} No finder verifies that a candidate contains one particle. TC merges
crossing tracks (the reason HDBSCAN was added, per its header); \code{SelectAndMergeTracks} attaches any
non-primary track with an end within 30~mm of a primary's far end; HDBSCAN's overlap join merges any two
same-direction clusters whose circles overlap by 25\,\%. A candidate that contains two particles is not
repaired by any ordering of its hits or clusters: z sorting, the nearest-neighbour walk and the arc walk
all produce one sequence through both. The greedy walk instead produces the U-turns measured by
\url{macro/Unpack_HDF5/a1975/D2_UKF/zuturn_dp.C}.

\paragraph{Cluster radius 0 on injected finders.} \code{AtPRA} initialises the cluster radius and distance
to 0. TC and Riemann pass them unchanged to \code{ClusterizeSmooth3D}, which then builds no clusters;
\code{SelectAndMergeTracks} rejects such tracks as degenerate and the fitters have no measurements. Always
call \code{SetClusterRadius} and \code{SetClusterDistance} on the algorithm before injecting it (the
\code{AtPRAtask} setters only apply on the legacy path).

\paragraph{Hits that vanish.} Hits of tracks removed by \code{SelectAndMergeTracks} or by the TC
\code{GeoRadius} filter are neither in a track nor in \code{GetNoiseHits()}. Do not compute
``assigned + noise = all hits'' efficiencies from TC output.

\paragraph{Pattern event header.} TC, HDBSCAN and Riemann create a fresh \code{AtPatternEvent}; the event
ID and the good/bad flag of the input \code{AtEvent} are not copied into it (\code{AtSampleConsensus}
copies them). Match events by entry number, or read the ID from the input event branch.

\paragraph{Legacy algorithm numbers 1 and 2.} They construct nothing; the first event inside the hit-number
window then dereferences an empty algorithm pointer.
